% chapters/16/evaluasi.tex
\bagian{Asesmen dan Instrumen Penilaian}
\label{sec:asesmen-playtest}

\begin{asesmenbox}
\textbf{Kedudukan asesmen bab ini.} Bab ini menuntaskan \SUBCPMK{4.3} pada
ranah C5, dan merupakan bab terakhir yang dinilai pada \CPMK{4}. Pengujian dan
presentasi pada bab ini termasuk bagian dari \textbf{Proyek Akhir}.

\vspace{0.4em}
Materi Bab~XVI termasuk bahan \textbf{Ujian Akhir Semester} pada minggu keenam
belas, dan seluruh keluarannya diserahkan sebagai \istilah{build},
\istilah{project source}, laporan pengembangan, serta lembar \istilah{playtest}
yang telah diisi penguji. Rubrik lengkap Proyek Akhir disajikan pada
Lampiran~A.
\end{asesmenbox}

\subsection*{Instrumen 1: Rubrik Bagian Pengujian dan Presentasi}

\begin{table}[htbp]
\centering
\caption{Rubrik penilaian pengujian dan presentasi beserta indikator yang diukur}
\label{tab:rubrik-praktikum-16}
{\setlength{\tabcolsep}{4pt}\footnotesize%
\begin{tabularx}{\textwidth}{@{}L{3.2cm} C{1.3cm} C{2.1cm} Y@{}}
\toprule
\headrow \textbf{Kriteria} & \textbf{Bobot} & \textbf{Indikator} & \textbf{Deskripsi capaian} \\
\midrule
\istilah{Checklist} fungsional & 25\% & \ind{a} &
Butirnya diturunkan dari GDD, mencakup ketiga kelompok, dapat diperiksa benar
atau salah, dan hasil setiap putaran tercatat. \\
\addlinespace
Pemeriksaan otomatis & 20\% & \ind{b} &
Level diperiksa dengan \code{ch16_uji_mandiri.py}; tempat yang terjebak
diperbaiki; jangkauan lompatan dihitung dari besaran levelnya sendiri. \\
\addlinespace
\istilah{Playtest} dan \istilah{log bug} & 30\% & \ind{a}, \ind{b} &
Lembar \istilah{playtest} terisi oleh penguji lain; setiap temuan dapat
diulang; berat setiap temuan ditimbang; temuan berat sudah diperbaiki. \\
\addlinespace
Laporan dan presentasi & 25\% & \ind{c} &
Demo disampaikan dalam waktu yang ditentukan; laporan memuat keadaan sebelum
dan sesudah; refleksi teknologi disertai contoh dari pekerjaan sendiri. \\
\bottomrule
\end{tabularx}}
\end{table}

\subsection*{Instrumen 2: Soal Konsep}

Jawablah dengan singkat. Kunci jawaban beserta indikator yang diukur disajikan
pada Tabel~\ref{tab:kunci-konsep-ch16}.

\begin{enumerate}
  \item Sebuah level memakai $v_0 = 640$ piksel per detik dan
        $g = 1600$ piksel per detik kuadrat, dengan petak berukuran 32 piksel.
        Hitunglah tinggi lompatan tertinggi dalam satuan petak.
  \item Jelaskan perbedaan \istilah{soft-lock} dan kegagalan program yang
        menghentikan permainan. Mengapa \istilah{soft-lock} lebih sulit
        ditemukan?
  \item Mengapa pemeriksaan yang memeriksa setiap tempat berdiri lebih dapat
        dipercaya daripada simulasi dengan titik awal acak?
  \item Sebuah catatan \istilah{bug} berbunyi ``musiknya berhenti sendiri''.
        Sebutkan dua keterangan yang perlu ditambahkan agar catatan itu dapat
        diulang.
  \item Kelompok Anda memutuskan untuk tidak mengerjakan sebuah mekanik yang
        sudah dijanjikan pada GDD. Apa yang sebaiknya dilakukan, dan mengapa?
\end{enumerate}

\begin{table}[htbp]
\centering
\caption{Kunci soal konsep Bab XVI beserta indikator yang diukur}
\label{tab:kunci-konsep-ch16}
{\setlength{\tabcolsep}{4pt}\footnotesize%
\begin{tabularx}{\textwidth}{@{}C{0.8cm} C{2.1cm} Y@{}}
\toprule
\headrow \textbf{No} & \textbf{Indikator} & \textbf{Pokok jawaban} \\
\midrule
1 & \ind{b} &
$h = 640^2 \div (2 \times 1600) = 409\,600 \div 3200 = 128$ piksel. Dalam
satuan petak, $128 \div 32 = 4{,}00$ petak. \\
\addlinespace
2 & \ind{b} &
\istilah{Soft-lock} membuat pemain tidak dapat melanjutkan walaupun programnya
masih berjalan dan tidak mencetak kesalahan apa pun. Ia lebih sulit ditemukan
sebab tidak ada tanda apa pun pada \istilah{Console}, dan hanya dialami oleh
pemain yang kebetulan berada di tempat yang salah. \\
\addlinespace
3 & \ind{b} &
Sebab pemeriksaan menyeluruh memeriksa seluruh keadaan, sedangkan simulasi
hanya menyentuh keadaan yang kebetulan terpilih. Kesalahan yang jarang muncul
dapat lolos sepenuhnya dari simulasi. Contohnya pada peta ketiga: simulasi
mengembalikan 82,0 persen, sedangkan pemeriksaan menyeluruh menemukan 15,6
persen tempat yang terjebak. \\
\addlinespace
4 & \ind{a}, \ind{b} &
Waktu kejadiannya, dan urutan tindakan yang mendahuluinya. Keterangan lain yang
berguna: tempat pemain berada, dan apakah kejadian itu dapat diulang. \\
\addlinespace
5 & \ind{c} &
Mekaniknya tetap dicantumkan pada GDD sebagai kekurangan yang diketahui,
disertai alasan mengapa tidak dikerjakan. Menghapusnya dari GDD agar hasilnya
tampak sesuai adalah ketidakjujuran akademik, dan sekaligus menghilangkan
bukti kemampuan menilai karya sendiri. \\
\bottomrule
\end{tabularx}}
\end{table}
